% Determinant by Laplace (cofactor) expansion along the first row: the textbook
% definition turned into an algorithm. Recursive, and deliberately so -- the
% point in class is that it costs O(n!) operations, versus O(n^3) for getting
% the determinant out of an LU factorization.
% Uses Chapra's notation: [A] for a matrix.
% Cropped to its own bounding box so it can be dropped onto a Xournal++ page as
% an image.
%
%   latexmk -pdf determinant-pseudocode.tex
\documentclass[border=6pt,varwidth=12cm]{standalone}

\usepackage{algorithm2e}
\usepackage{amsmath,amssymb}

\RestyleAlgo{plain}          % no float, no ruled box
\LinesNumbered
\DontPrintSemicolon
\SetKwComment{Comment}{// }{}
\SetKwInOut{Input}{input}
\SetKwInOut{Output}{output}
\SetKwFunction{Det}{Det}
\SetKwProg{Fn}{function}{}{}

\newcommand{\mat}[1]{[#1]}       % matrix

\begin{document}
\begin{minipage}{\linewidth}
\textbf{Determinant by Laplace Expansion}\par\medskip
\begin{algorithm}[H]
  \Input{square matrix $\mat{A}$}
  \Output{$\det \mat{A}$}
  \BlankLine
  \Fn{\Det{$\mat{A}$}}{
    $n \gets \text{nrows}(\mat{A})$\;
    \If{$n = 1$}{
      \Return $a_{11}$
    }
    $d \leftarrow 0$\;
    \For{$j \leftarrow 1$ \KwTo $n$}{
        $\mat{M} \gets \mat{A}$ with row $1$ and column $j$ deleted\;
        $d \leftarrow d + (-1)^{1+j}\, a_{1j} \cdot \Det{$\mat{M}$}$\;
    }
    \Return $d$\;
  }
\end{algorithm}
\end{minipage}
\end{document}
